% !TEX root = ../../ctfp-print.tex

\lettrine[lhang=0.17]{类}{型与函数范畴}（the category of types and functions）在编程中扮演着重要角色，因此让我们先讨论类型的本质及其必要性。

\section{谁需要类型？}

关于静态类型（static）与动态类型（dynamic）、强类型（strong）与弱类型（weak）系统的优劣似乎存在诸多争议。我将通过一个思想实验来阐释这些选择的本质区别：设想有数百万只猴子在电脑键盘前欢快地随机敲击按键，生成程序并进行编译运行。

\begin{figure}[H]
  \centering
  \includegraphics[width=0.3\textwidth]{images/img_1329.jpg}
\end{figure}

\noindent
在机器语言层面，猴子产生的任意字节组合都会被接受并执行。但在高级语言中，我们更期待编译器能够检测词汇和语法错误。许多猴子将因此得不到香蕉奖励，但剩余程序的可用性会显著提升。类型检查则为荒谬程序设置了又一道防线。此外，在动态类型语言中，类型不匹配会在运行时暴露；而在静态强类型语言中，这种错误会被提前至编译阶段发现，从而在程序运行前就消除大量潜在错误。

那么问题在于：我们究竟想取悦猴子，还是生产正确的程序？

"打字猴子"思想实验的经典目标是产出莎士比亚全集。此时引入拼写检查和语法检查已能大幅提高成功率。而类型检查的类比机制将进一步确保：一旦罗密欧（Romeo）被声明为人类，他就不会突然长出树叶，或在其强大引力场中捕获光子。

\section{类型关乎可组合性}

范畴论（category theory）的核心在于箭头的组合。但并非任意两个箭头都能组合——前一个箭头的目标对象必须与后一个箭头的源对象相同。在编程中，我们常将一个函数的结果传递给另一个函数。若目标函数无法正确解析源函数产生的数据，程序就会失效。这种情况下两端必须匹配才能完成组合，而语言的类型系统越强大，就越能精确描述并机械验证这种匹配。

我听到的唯一有力的反对方观点认为强静态类型检查可能会排除某些语义正确的程序。实际上这种情况极少发生，而且任何语言都会提供某种绕过类型系统的后门机制（backdoor）。即便是Haskell也包含\code{unsafeCoerce}这样的非安全强制转换函数。但这类机制应当审慎使用。正如卡夫卡（Franz Kafka）笔下的格里高尔·萨姆沙（Gregor Samsa）变成甲虫时突破了类型系统，我们都知道故事的结局如何。

另一个常见观点是类型系统会给程序员带来过多负担。当我不得不手动编写C++迭代器声明时确实对此深有体会，不过存在一种名为\newterm{类型推断}（type inference）的技术，它允许编译器通过上下文自动推导出大部分类型信息。如今C++已支持通过\code{auto}关键字声明变量，让编译器自动判断其类型。

在Haskell中，除特殊情况外类型注解完全是可选的。但程序员仍倾向于使用类型注解，因为它们能揭示代码的语义特征，并使编译错误更易理解。Haskell社区的普遍实践是在项目初期先设计类型系统，随后这些类型注解将驱动具体实现，并成为编译器强制校验的代码注释。\sloppy{}

强静态类型常被当作不进行测试的借口。你或许听过Haskell程序员说"只要能编译通过，必定正确"。当然，这并不能保证类型正确的程序一定能产生预期输出。这种轻率态度导致多项研究显示Haskell在代码质量上的优势并不如预期显著。商业环境下，修复bug的压力往往只维持到某个质量阈值，这更多与软件开发成本和终端用户容忍度相关，而非编程语言本身或开发方法论。更好的评估标准应是统计延期交付或功能严重缩水的项目数量。

至于"单元测试可替代强类型系统"的观点，考虑强类型语言中常见的重构实践：修改某个函数参数的类型。在强类型语言中只需修改函数声明并修复所有编译错误即可。而在弱类型语言中，参数类型的变更无法传播到调用点。虽然单元测试可能捕获部分错误，但测试本质上是概率性而非确定性的过程。测试永远无法替代形式化证明。

\section{类型是什么？}

类型的最简单直觉理解是它们代表值的集合。类型 \code{Bool}（记住，具体类型在Haskell中以大写字母开头）是由 \code{True} 和 \code{False} 组成的二元集合。类型 \code{Char} 是包含所有Unicode字符（如 \code{a} 或 \code{ą}）的集合。

集合可以是有限或无限的。\code{String} 类型（本质上是 \code{Char} 列表的同义词）就是一个无限集合的示例。

当我们声明 \code{x} 为 \code{Integer} 类型时：

\src{snippet01}
我们是在说它是整数集合中的一个元素。Haskell中的 \code{Integer} 是无限集合，支持任意精度算术运算。此外还有一个对应机器类型的有限集合 \code{Int}，类似于C++的 \code{int}。

这种将类型与集合等同的观念存在一些细微差别。当涉及多态函数的循环定义时会出现问题，而且你无法构造包含所有集合的"集合"；不过正如我承诺过的，我不会过分拘泥于数学细节。关键在于确实存在一个名为 $\Set$ 的集合范畴，我们将以此为基础进行讨论。在 $\Set$ 中，对象是集合，态射（箭头）是函数。

$\Set$ 是一个非常特殊的范畴，因为我们能实际观察其对象的内部结构并获得大量直观认识。例如我们知道空集没有任何元素，知道存在特殊单元素集合，知道函数将一个集合的元素映射到另一个集合的元素——可以将两个元素映射为一个，但不能将一个元素映射为两个。我们知道恒等函数将每个元素映射到自身等等。我们的计划是逐步遗忘这些具体信息，转而用纯粹的范畴论语言（即通过对象和箭头）重新表述这些概念。

在理想情况下，我们可以简单地说Haskell类型就是集合，Haskell函数就是集合间的数学函数。但存在一个关键问题：数学函数不执行任何计算代码——它直接给出答案，而Haskell函数必须通过计算得到结果。如果答案能在有限步骤内获得（无论这个数目有多大），这都不是问题。但某些递归计算可能导致程序永不终止。我们不能简单地从Haskell中禁止非终止函数，因为判定函数是否终止是不可判定问题（著名的停机问题）。这就是计算机科学家提出的一个绝妙想法（或者说是重大hack，取决于你的立场）：通过扩展每个类型增加一个特殊值称为\newterm{底值(bottom)}，记作 \code{\_|\_} 或Unicode符号 $\bot$。这个"值"对应非终止计算。因此声明为：

\src{snippet02}
的函数可能返回 \code{True}、\code{False} 或 \code{\_|\_}（表示永不终止）。

有趣的是，一旦接受底值作为类型系统的一部分，就可以方便地将所有运行时错误视为底值，甚至允许函数显式返回底值。后者通常通过表达式 \code{undefined} 实现，例如：

\src{snippet03}
该定义能通过类型检查，因为 \code{undefined} 求值到底值，而底值属于任何类型（包括 \code{Bool}）。你甚至可以写：

\src{snippet04}
（不带 \code{x}），因为底值同样属于类型 \code{Bool -> Bool}。

可能返回底值的函数被称为偏函数（partial function），与之相对的全函数（total function）则对所有可能参数都返回有效结果。

由于底值的存在，你会看到Haskell类型与函数构成的范畴被称为 $\Hask$ 而非 $\Set$。从理论角度看，这是无穷麻烦的根源，所以此时我将用我的"屠夫刀"果断切断这条推理路线。从实用主义角度看，我们可以忽略非终止函数和底值，将 $\Hask$ 视为真正的 $\Set$.\footnote{Nils Anders Danielsson, John Hughes, Patrik Jansson, Jeremy Gibbons, \href{http://www.cs.ox.ac.uk/jeremy.gibbons/publications/fast+loose.pdf}{Fast and Loose Reasoning is Morally Correct}. 该论文论证了在多数场景下忽略底值的合理性。}

\section{为什么需要一个数学模型？}

作为一名程序员，你对编程语言的语法（syntax）和文法（grammar）非常熟悉。语言的这些方面通常在语言规范的开篇就用形式化记号进行描述。然而语言的意义（meaning），即其语义（semantics），却难以精确描述：它往往需要更多篇幅，很少足够形式化，几乎从未完整覆盖全部细节。因此才会产生语言律师们永无止境的争论，以及大量专注于解析语言标准细微之处的著作。

我们确实有用于描述语言语义的形式化工具，但由于其复杂性，这些工具大多应用于简化后的教学语言，而非现实世界中的编程巨兽。其中一种名为\newterm{操作语义学}(operational semantics)的工具描述程序执行的机械过程，它定义了一个形式化的理想化解释器。工业级语言（如C++）的语义通常通过非形式化的操作推理进行描述，这种描述往往基于"抽象机器"的概念。

问题在于使用操作语义学证明程序性质极其困难。要展示某个程序的性质，本质上需要通过理想化解释器"运行"该程序。

程序员从不进行形式化正确性证明的事实无关紧要。我们总是"认为"自己编写的是正确程序。没有人会坐在键盘前说："哦，我就随便写几行代码看看结果"。我们相信所写的代码会执行特定操作并产生预期结果，当结果不符时通常会感到惊讶。这表明我们确实会对编写程序进行推理，只是通过大脑中的解释器进行模拟。但人类很难跟踪所有变量变化——计算机擅长运行程序，人类却不擅长！若非如此，我们就不需要计算机了。

但存在另一种选择：\newterm{指称语义学}(denotational semantics)，它基于数学理论。在指称语义学中，每个编程结构都有对应的数学解释。通过这种方式，若要证明程序性质，只需证明相应的数学定理即可。虽然定理证明看似困难，但人类已积累数千年数学方法，拥有丰富的知识储备可供调用。而且相比专业数学家证明的定理，编程中遇到的问题通常相当简单（甚至微不足道）。

考虑Haskell语言中阶乘函数的定义（该语言特别适合指称语义学）：

\src{snippet05}
表达式\code{{[}1..n{]}表示从\code{1}到\code{n}的整数列表，\code{product}函数将列表所有元素相乘。这就像数学教材中的阶乘定义。对比C语言实现：

\begin{snip}{c}
int fact(int n) {
    int i;
    int result = 1;
    for (i = 2; i <= n; ++i)
        result *= i;
    return result;
}
\end{snip}
还需要进一步比较吗？

当然，我承认这是个简单的例子！阶乘函数具有明显的数学指称。敏锐的读者可能会问：从键盘读取字符或通过网络发送数据包的数学模型是什么？很长时间以来，这个问题会导致复杂的解释。看起来指称语义学难以胜任许多重要任务（这些任务对编写实用程序至关重要，却能被操作语义学轻松处理）。突破来自范畴论：Eugenio Moggi发现计算效应可映射到单子(monads)。这一发现不仅让指称语义学重获新生，提升了纯函数式程序的可用性，还为传统编程提供了新视角。当我们掌握更多范畴论工具后，将进一步讨论单子概念。

编程数学模型的重要优势之一是支持软件正确性的形式化证明。当你开发消费类软件时这似乎并不重要，但在某些领域，程序失败的代价可能极高甚至危及生命。即使开发医疗系统Web应用时，你也会感激Haskell标准库中的函数和算法附带正确性证明这一事实。

\section{纯函数与不纯函数}

我们在C++或任何其他命令式语言中称为"函数"(function)的概念，
与数学家所定义的函数并不完全相同。数学中的函数本质上只是值到值的映射关系。

我们可以在编程语言中实现数学意义上的函数：这种函数接受输入值后会计算输出值。
实现数字平方的函数很可能通过将输入值自乘来完成运算。每次调用该函数时都会执行同样的计算，
且保证相同输入必定产生相同输出。数字的平方不会随着月相周期发生变化。

此外，计算数值平方的操作不应该产生副作用（side effect），例如给你的狗分发美味零食。
如果某个"函数"具有此类副作用，则难以用数学函数进行建模。

在编程语言中，那些对于相同输入始终产生相同结果且没有副作用的函数被称为\newterm{纯函数}(pure function)。
在Haskell这类纯函数式语言中，所有函数都是纯函数。正因如此，这类语言更容易赋予指称语义(denotational semantics)，
并使用范畴论(category theory)进行建模。对于其他类型的语言，我们总可以通过自我约束的方式
仅使用纯函数子集，或将副作用单独进行分析。稍后我们将看到单子(monad)如何让我们仅通过纯函数
就能模拟各种类型的副作用。因此，通过限制自己使用数学意义上的函数并不会造成表达能力的损失。

\section{类型的例子}

一旦你意识到类型就是集合，就可以思考一些相当奇特的类型。例如，空集对应的类型是什么？不，它不是C++的\code{void}，虽然这个类型在Haskell中确实被称为\code{Void}。这是一个没有任何值 inhabit 的类型。你可以定义一个接受\code{Void}的函数，但你永远无法调用它。要调用它，你需要提供一个\code{Void}类型的值，而这样的值根本不存在。至于这个函数能返回什么，则完全没有限制。它可以返回任何类型（尽管它永远不会返回，因为它无法被调用）。换句话说，这是一个返回类型多态的函数。Haskellers给它起了个名字：

\src{snippet06}
（记住，\code{a}是一个类型变量，可以代表任何类型。）这个名字并非偶然。在逻辑领域有一个更深层的类型和函数解释，称为Curry-Howard同构（Curry-Howard isomorphism）。类型\code{Void}表示假命题，而函数\code{absurd}的类型对应着"从假命题可以推出任何结论"这一逻辑陈述，正如拉丁谚语"ex falso sequitur quodlibet"所表达的那样。

接下来是单元素集合对应的类型。这种类型只有一个可能的值。这个值就"存在在那里"。你可能不会立即认识到这一点，但这就是C++的\code{void}。考虑从此类型到此类型的函数：一个来自\code{void}的函数总是可以被调用的。如果是纯函数，它总会返回相同的结果。以下是此类函数的示例：

\begin{snip}{c}
int f44() { return 44; }
\end{snip}
你可能会认为这个函数接收"nothing"，但正如我们刚才看到的，一个接收"nothing"的函数永远无法被调用，因为没有值能表示"nothing"。那么这个函数实际上接收什么呢？概念上，它接收一个只有单一实例的虚拟值，因此我们无需显式提及它。然而在Haskell中，这个值有专门的符号：一对空括号\code{()}。所以，有趣的是（或者这真的是巧合吗？），在C++和Haskell中调用无参函数的形式看起来是一样的。此外，由于Haskell对简洁性的追求，同一个符号\code{()}被用于表示类型、构造器以及对应单元素集合的唯一值。这是Haskell中的这个函数：

\src{snippet07}
第一行声明\code{f44}将类型\code{()}（发音为"unit"）映射到类型\code{Integer}。第二行通过模式匹配unit的唯一构造器\code{()}来定义\code{f44}，并返回数字44。你可以通过提供unit值\code{()}来调用这个函数：

\begin{snip}{c}
f44 ()
\end{snip}
请注意，每个unit函数都等价于从目标类型中选择一个单一元素。事实上，你可以将\code{f44}视为数字44的不同表示形式。这展示了如何用函数（箭头）代替显式的集合元素表示法。从unit到任意类型$A$的函数与该集合$A$的元素之间存在一一对应关系。

关于返回\code{void}类型的函数（在Haskell中即unit返回类型）呢？在C++中这类函数用于产生副作用，但我们知道这些并不是数学意义上的真正函数。返回unit的纯函数不做任何事情：它会丢弃其参数。

数学上，从集合$A$到单元素集合的函数会将$A$的每个元素映射到该单元素集合的唯一元素。对于每个$A$，恰好存在这样一个函数。这是针对\code{Integer}的实现：

\src{snippet08}
你给它任何整数，它都会返回unit。为了简洁起见，Haskell允许你使用下划线作为通配符模式来表示被丢弃的参数。这样你就不用为它发明一个名字了。因此上面的代码可以重写为：

\src{snippet09}
请注意，这个函数的实现不仅不依赖于传入的值，甚至不依赖于参数的类型。

可以用相同公式实现任意类型的函数被称为参数化多态（parametric polymorphism）。你可以用一个方程配合类型参数而非具体类型来实现这一系列函数。我们应该如何命名从任意类型到unit类型的多态函数？当然称之为\code{unit}：

\src{snippet10}
在C++中你会这样编写这个函数：

\begin{snip}{cpp}
template<class T>
void unit(T) {}
\end{snip}
在类型学分类中下一个类型是二元集合。在C++中称为\code{bool}，在Haskell中理所当然地称为\code{Bool}。区别在于C++的\code{bool}是内建类型，而在Haskell中可以如下定义：

\src{snippet11}
（这个定义的解读方式是：\code{Bool}要么是\code{True}，要么是\code{False}。）原则上，也应该能够在C++中将布尔类型定义为枚举：

\begin{snip}{cpp}
enum bool {
    true,
    false
};
\end{snip}
但C++的\code{enum}本质上是整数。虽然可以改用C++11的``\code{enum class}''，但这会导致必须用类名限定其值，如\code{bool::true}和\code{bool::false}，更不用说每个使用它的文件都需要包含相应头文件。

从\code{Bool}出发的纯函数只是从目标类型中选取两个值，一个对应\code{True}，另一个对应\code{False}。

返回\code{Bool}的函数被称为\newterm{谓词（predicate）}。例如，Haskell的\code{Data.Char}库充满了像\code{isAlpha}或\code{isDigit}这样的谓词。在C++中有类似的库\code{}定义了\code{isalpha}和\code{isdigit}，但它们返回的是\code{int}而非布尔值。真正的谓词定义在\code{std::ctype}中，形式为\code{ctype::is(alpha, c)}，\code{ctype::is(digit, c)}等。

\section{挑战}

\begin{enumerate}
  \tightlist
  \item
        在你最喜欢的编程语言中定义一个高阶函数（或函数对象）\code{memoize}。
        该函数接受一个纯函数\code{f}作为参数，并返回一个行为与\code{f}几乎完全相同的函数，
        只不过它对每个参数只调用原始函数一次，内部存储计算结果，并在后续相同参数调用时直接返回存储结果。
        通过观察性能可以区分记忆化函数与原始函数。例如，尝试记忆化一个计算耗时较长的函数：
        首次调用时需要等待计算结果，但后续相同参数调用应立即获得结果。
  \item
        尝试对你常用的标准库中生成随机数的函数进行记忆化。这种方法有效吗？
  \item
        大多数随机数生成器都可以通过种子初始化。实现一个接受种子参数、
        使用该种子调用随机数生成器并返回结果的函数。对该函数进行记忆化，效果如何？
  \item
        下列C++函数中哪些是纯函数？尝试对它们进行记忆化并观察多次调用时的记忆化效果与非记忆化表现。

        \begin{enumerate}
          \tightlist
          \item
                文中示例中的阶乘函数。
          \item
                \begin{minted}{cpp}
std::getchar()
\end{minted}
          \item
                \begin{minted}{cpp}
bool f() {
    std::cout << "Hello!" << std::endl;
    return true;
}
\end{minted}
          \item
                \begin{minted}{cpp}
int f(int x) {
    static int y = 0;
    y += x;
    return y;
}
\end{minted}
        \end{enumerate}
  \item
        从\code{Bool}到\code{Bool}的不同函数共有多少种？你能全部实现吗？
  \item
        绘制一个范畴图，其唯一对象为类型\code{Void}、\code{()}（单位类型）和\code{Bool}；
        箭头对应这些类型间的所有可能函数，并用函数名称标记各箭头。
\end{enumerate}